\section{Conclusion}

La recherche exacte confirme la théorie : la liste est en $O(N)$, les deux
arbres en $O(\log N)$. À $N=10^6$, la liste prend 39{,}9 s contre moins de
0{,}001 s pour le BST 2D.

Le résultat le plus surprenant est que le BST Morton devient plus lent
que la simple liste dès que le rayon augmente.
Ce comportement s'explique par deux effets combinés.

\paragraph{1. Faux positifs}
Pour chercher des points, on calcule les codes Z des deux coins d'un carré.
Ensuite, on cherche dans l'arbre tous les éléments entre ces deux codes.
Le problème est que la courbe en Z passe par de grandes zones qui n'ont rien à voir avec notre carré.
Résultat : on doit vérifier beaucoup plus de points que nécessaire.

\paragraph{2. Gestion de la mémoire}
Avec les codes de Morton, l'algorithme crée un nœud en mémoire pour chaque élément trouvé, même s'il est hors du rayon.
Ensuite, il doit supprimer ces nœuds inutiles. C'est très lourd.
Au contraire, la recherche linéaire n'utilise de la mémoire que pour les bons éléments.
Cette perte de temps pour créer et supprimer des éléments inutiles ralentit beaucoup l'algorithme.

\paragraph{Conclusion}

L'arbre de recherche 2D reste la solution la plus rapide.
Il permet de "couper" des parties entières du plan (pruning), ce qui est efficace même pour de grands rayons.
C’est donc la meilleure solution pour une application de taxi à Porto.
